\documentclass[../distribution_theory_notes.tex]{subfiles}
\begin{document}
\section{Aula 27 - 29 de Novembro, 2024}
\subsection{Transformada de Fourier no Espaço de Schwartz}
Na aula passada estudamos as propriedades da transformada de Fourier de funções integráveis. Para $f\in L^1(\mathbb R^n)$,
\[
 \widehat f(\xi)=\int_{\mathbb R^n}e^{-2\pi i x\cdot\xi}f(x)\,dx,
\]
e o operador $\mathcal F$ satisfaz, entre outras, as propriedades:
\begin{enumerate}
 \item
 \[
  \mathcal F:L^1\longrightarrow C_0,
  \qquad
  \|\widehat f\|_\infty\leq\|f\|_{L^1};
 \]
 \item
 \[
  \widehat{\tau_h f}(\xi)
  =e^{-2\pi i h\cdot\xi}\widehat f(\xi);
 \]
 \item se $T:\mathbb R^n\to\mathbb R^n$ é um isomorfismo,
 \[
  \widehat{f\circ T}(\xi)
  =|\det T|^{-1}\widehat f((T^*)^{-1}\xi);
 \]
 \item
 \[
  \widehat{f*g}=\widehat f\,\widehat g;
 \]
 \item
 \[
  \int\widehat f\,g=\int f\,\widehat g;
 \]
 \item
 \[
  \partial_\xi^\alpha\widehat f
  =\left[(-2\pi i x)^\alpha f\right]^\wedge;
 \]
 \item sob hipóteses naturais de integrabilidade e decaimento,
 \[
  \widehat{\partial^\alpha f}
  =(2\pi i\xi)^\alpha\widehat f.
 \]
\end{enumerate}
Recordemos também duas transformadas explícitas:
\[
 \widehat{\chi_{[-a,a]}}(\xi)=\frac{\sin(2\pi a\xi)}{\pi\xi},
\]
com o valor $2a$ em $\xi=0$, e
\[
 \widehat{e^{-\pi a|x|^2}}(\xi)
 =a^{-n/2}e^{-\pi|\xi|^2/a},
 \qquad a>0.
\]
Em particular, $e^{-\pi|x|^2}$ é um ponto fixo de $\mathcal F$.

Queremos agora analisar a transformada no espaço de Schwartz $\mathcal S(\mathbb R^n)$. O resultado central é que $\mathcal F$ é um isomorfismo de $\mathcal S$ em $\mathcal S$, permitindo estendê-la por dualidade ao espaço $\mathcal S'$ das distribuições temperadas.

\subsection{Seminormas do Espaço de Schwartz}
Para $f\in\mathcal C^\infty(\mathbb R^n)$, consideremos
\[
 p_{(\alpha,\beta)}(f)
 =\sup_{x\in\mathbb R^n}|x^\alpha\partial^\beta f(x)|
\]
e
\[
 q_{(\alpha,\beta)}(f)
 =\sup_{x\in\mathbb R^n}|\partial^\beta(x^\alpha f)(x)|.
\]

\begin{lemma*}
 As seguintes condições são equivalentes:
 \begin{enumerate}
  \item $f\in\mathcal S$;
  \item $p_{(\alpha,\beta)}(f)<\infty$ para todos $\alpha,\beta$;
  \item $q_{(\alpha,\beta)}(f)<\infty$ para todos $\alpha,\beta$.
 \end{enumerate}
 Além disso, as famílias $(p_{(\alpha,\beta)})$ e $(q_{(\alpha,\beta)})$ definem a mesma topologia em $\mathcal S$.
\end{lemma*}

\begin{proof*}
 Pela regra de Leibniz,
 \[
  \partial^\beta(x^\alpha f)
  =\sum_{\gamma\leq\beta}
  \binom{\beta}{\gamma}
  \partial^\gamma(x^\alpha)\,
  \partial^{\beta-\gamma}f.
 \]
 Cada $\partial^\gamma(x^\alpha)$ é um múltiplo de um monômio ou é zero. Assim, cada $q_{(\alpha,\beta)}$ é majorada por uma combinação linear finita de seminormas $p_{(\mu,\nu)}$. Isolando o termo $x^\alpha\partial^\beta f$ na mesma identidade e procedendo por indução em $|\beta|$, obtemos a estimativa inversa. \qedsymbol
\end{proof*}

\begin{lemma*}
 Para todo $1\leq p\leq\infty$,
 \[
  \mathcal S(\mathbb R^n)\hookrightarrow L^p(\mathbb R^n)
 \]
 continuamente.
\end{lemma*}

\begin{proof*}
 O caso $p=\infty$ é imediato. Para $p<\infty$, escolha $N\in\mathbb Z_+$ com $Np>n$. Então
 \begin{align*}
  \int_{\mathbb R^n}|f(x)|^p\,dx
  &=\int(1+|x|)^{-Np}
  \bigl[(1+|x|)^N|f(x)|\bigr]^p\,dx\\
  &\leq C\|f\|_{(N,0)}^p
  \int_{\mathbb R^n}(1+|x|)^{-Np}\,dx<\infty.
 \end{align*}
 A última integral é finita por coordenadas polares. \qedsymbol
\end{proof*}

\begin{theorem*}
 A transformada de Fourier
 \[
  \mathcal F:\mathcal S\longrightarrow\mathcal S
 \]
 está bem definida e é contínua.
\end{theorem*}

\begin{proof*}
 Seja $f\in\mathcal S$. Como, para todos $\alpha,\beta$, $x^\alpha\partial^\beta f\in L^1$, podemos usar as propriedades de derivação e multiplicação. Assim,
 \begin{align*}
  \partial_\xi^\beta(\xi^\alpha\widehat f)(\xi)
  &=(2\pi i)^{-|\alpha|}
  \partial_\xi^\beta\left[\widehat{\partial_x^\alpha f}\right](\xi)\\
  &=(2\pi i)^{-|\alpha|}
  \left[(-2\pi i x)^\beta\partial_x^\alpha f\right]^\wedge(\xi).
 \end{align*}
 Portanto
 \[
  q_{(\alpha,\beta)}(\widehat f)
  \leq(2\pi)^{|\beta|-|\alpha|}
  \|x^\beta\partial^\alpha f\|_{L^1}. \tag{*}
 \]

 Escolha $N>n+|\beta|$. Então
 \begin{align*}
  \|x^\beta\partial^\alpha f\|_{L^1}
  &\leq\int(1+|x|)^{-N}
  (1+|x|)^N|x|^{|\beta|}|\partial^\alpha f(x)|\,dx\\
  &\leq C_N\|f\|_{(N+|\beta|,\alpha)}.
 \end{align*}
 Substituindo em (*),
 \[
  q_{(\alpha,\beta)}(\widehat f)
  \leq C_{N,\alpha,\beta}
  \|f\|_{(N+|\beta|,\alpha)}.
 \]
 Isso mostra que $\widehat f\in\mathcal S$ e, ao mesmo tempo, a continuidade de $\mathcal F$. \qedsymbol
\end{proof*}

\begin{tcolorbox}[
 skin=enhanced,
 title=Observação,
 fonttitle=\bfseries,
 colframe=black,
 colbacktitle=cyan!75!white,
 colback=cyan!15,
 colbacklower=black,
 coltitle=black,
 drop fuzzy shadow]
 O teorema anterior fornece um espaço invariante pela transformada de Fourier. Isso é essencial para estender $\mathcal F$ por dualidade. A mesma estratégia não funciona diretamente em $\mathcal D'$ porque uma função teste não permanece, em geral, com suporte compacto após a transformada; isso será formalizado pelo teorema de Paley--Wiener.
\end{tcolorbox}

\subsection{Transformada de Fourier de Distribuições Temperadas}
Denotemos por $\mathcal S'=\mathcal S'(\mathbb R^n)$ o dual topológico de $\mathcal S$. Para $u\in\mathcal S'$, definimos sua transformada de Fourier $\widehat u=\mathcal F(u)$ por
\[
 \langle\widehat u,\varphi\rangle
 :=\langle u,\widehat\varphi\rangle,
 \qquad \forall\varphi\in\mathcal S.
\]
Essa definição faz sentido porque $\widehat\varphi\in\mathcal S$ e $\mathcal F:\mathcal S\to\mathcal S$ é contínua. Em particular, $\widehat u$ é linear e contínua em $\mathcal S$.

Se $\mathcal S'$ é munido da topologia fraca $*$, a transformada é contínua, pois, para cada $\varphi\in\mathcal S$,
\[
 p_\varphi(\widehat u)
 =|\langle\widehat u,\varphi\rangle|
 =|\langle u,\widehat\varphi\rangle|
 =p_{\widehat\varphi}(u).
\]

Também é útil recordar a reflexão
\[
 (Rf)(x)=f(-x).
\]
Com nossas convenções,
\[
 \mathcal F^{-1}=R\circ\mathcal F
\]
no nível das funções para as quais a inversão já é válida.

\begin{tcolorbox}[
 skin=enhanced,
 title=Exercício,
 fonttitle=\bfseries,
 colframe=black,
 colbacktitle=cyan!75!white,
 colback=cyan!15,
 colbacklower=black,
 coltitle=black,
 drop fuzzy shadow]
 Mostre que $\mathcal E'\hookrightarrow\mathcal S'$. Uma maneira é usar que $\mathcal C^\infty\supset\mathcal S$ e que as estimativas que definem uma distribuição de suporte compacto são controladas pelas seminormas de Schwartz.
\end{tcolorbox}

\subsection{Inversão da Transformada de Fourier}
Usaremos os fatos seguintes.
\begin{enumerate}
 \item Para $a>0$,
 \[
  f_a(x)=e^{-\pi a|x|^2}
  \quad\Longrightarrow\quad
  \widehat f_a(\xi)=a^{-n/2}e^{-\pi|\xi|^2/a}.
 \]
 \item Tomando $a=4\pi t$, $t>0$,
 \[
  \widehat f_{4\pi t}(\xi)
  =\frac{1}{(4\pi t)^{n/2}}e^{-|\xi|^2/(4t)}.
 \]
 \item Definindo
 \[
  g(\xi)=\pi^{-n/2}e^{-|\xi|^2},
  \qquad
  g_s(\xi)=s^{-n}g(\xi/s),
 \]
 temos $\int g=1$ e $(g_s)$ é uma família regularizante.
 \item Se $f\in L^1$, então
 \[
  g_s*f\longrightarrow f\quad\text{em }L^1,
  \qquad s\to0^+.
 \]
 \item Para $f,g\in L^1$,
 \[
  \int\widehat f\,g=\int f\,\widehat g.
 \]
\end{enumerate}

\begin{theorem*}[Inversão da Transformada de Fourier]
 Seja $f\in L^1(\mathbb R^n)$ tal que $\widehat f\in L^1(\mathbb R^n)$. Então $f$ coincide quase em todo ponto com uma função contínua $f_0$ dada por
 \[
  f_0(x)=\int_{\mathbb R^n}e^{2\pi i x\cdot\xi}\widehat f(\xi)\,d\xi.
 \]
 Em outras palavras,
 \[
  f=\check{\widehat f}\qquad\text{q.t.p.}
 \]
\end{theorem*}

\begin{proof*}
 Para $s>0$, $g_s*f$ é contínua e
 \[
  g_s*f\longrightarrow f\quad\text{em }L^1.
 \]
 Por outro lado, usando a identidade de dualidade da transformada e a transformada explícita da gaussiana, obtemos
 \begin{align*}
  (g_s*f)(x)
  &=\int_{\mathbb R^n}
  e^{-c_s|\xi|^2}e^{2\pi i x\cdot\xi}\widehat f(\xi)\,d\xi,
 \end{align*}
 onde $c_s\to0$ quando $s\to0$. Como $\widehat f\in L^1$ e o fator gaussiano tem módulo menor ou igual a $1$, pelo teorema da convergência dominada,
 \[
  (g_s*f)(x)
  \longrightarrow
  \int e^{2\pi i x\cdot\xi}\widehat f(\xi)\,d\xi
  =:f_0(x).
 \]
 A função $f_0$ é contínua pelo mesmo argumento usado para provar a continuidade da transformada de uma função $L^1$.

 Como $g_s*f\to f$ em $L^1$, existe uma subsequência convergindo para $f$ quase em todo ponto. O limite pontual dessa subsequência também é $f_0$, portanto $f=f_0$ q.t.p. \qedsymbol
\end{proof*}

\begin{crl*}
 A transformada de Fourier é um isomorfismo de $\mathcal S$ em $\mathcal S$.
\end{crl*}

\begin{proof*}
 Se $f\in\mathcal S$, então $f,\widehat f\in L^1$, e o teorema de inversão se aplica. Assim,
 \[
  \mathcal F^{-1}=R\circ\mathcal F
 \]
 em $\mathcal S$. Como $R$ e $\mathcal F$ são contínuos em $\mathcal S$, a inversa também é contínua. \qedsymbol
\end{proof*}

\begin{tcolorbox}[
 skin=enhanced,
 title=Exemplo,
 fonttitle=\bfseries,
 colframe=black,
 colbacktitle=cyan!75!white,
 colback=cyan!15,
 colbacklower=black,
 coltitle=black,
 drop fuzzy shadow]
 Para todo $1\leq p\leq\infty$,
 \[
  L^p(\mathbb R^n)\hookrightarrow\mathcal S'(\mathbb R^n).
 \]
 Se $1\leq p<\infty$ e $p'$ é o expoente conjugado, então $\mathcal S\hookrightarrow L^{p'}$; para $f\in L^p$,
 \[
  \left|\int f(x)\varphi(x)\,dx\right|
  \leq\|f\|_{L^p}\|\varphi\|_{L^{p'}},
 \]
 que é controlado por uma seminorma de Schwartz. O caso $p=\infty$ é análogo, usando $\mathcal S\hookrightarrow L^1$.
\end{tcolorbox}

\end{document}
